\documentclass[10pt,a4paper]{amsart}
\usepackage[margin=19mm]{geometry}
\usepackage{amsmath,amssymb,mathtools}
\usepackage[scr=rsfs,cal=euler]{mathalfa}
\usepackage{xfrac}
\usepackage[unicode,hidelinks,bookmarks=false]{hyperref}
\newcommand{\ie}{i.e.\ }
\newcommand{\cf}{cf.\ }
\newcommand{\resp}{respectively\ }
\newcommand{\Subsection}{\subsection}
\providecommand{\nobreakdash}{\nobreak\hbox{-}}
\setlength{\emergencystretch}{2em}
\multlinegap0pt
\usepackage{tikz}
\usetikzlibrary{cd}
\usepackage{amssymb}
\usepackage{mathtools}
\usepackage{dsfont}
\newtheorem{lemma}{Lemma}
\newtheorem{remark}{Remark}
\newcommand{\R}{{\mathbb R}}
\newcommand{\T}{{\mathbb T}}
\newcommand{\cA}{{\mathcal A}}
\newcommand{\cB}{{\mathcal B}}
\newcommand{\cC}{{\mathcal C}}
\newcommand{\cD}{{\mathcal D}}
\newcommand{\cI}{{\mathcal I}}
\newcommand{\cK}{{\mathcal K}}
\newcommand{\cV}{{\mathcal V}}
\newcommand{\cX}{{\mathcal X}}
\renewcommand{\emptyset}{\varnothing}
\DeclareMathOperator{\interior}{int}
\DeclareMathOperator*{\toup}{\longrightarrow} 
\newcommand{\twod}{{2d}}
\newcommand{\vvv}{|\kern-0.25ex|\kern-0.25ex|}
\title{\bf Equivariant critical point theory and bifurcation of $3d$ gravity-capillary Stokes waves}
\author{Tommaso Barbieri, Massimiliano Berti, Marco Mazzucchelli}
\date{Source: arXiv:2603.27847v2 (CC BY 4.0)}
\begin{document}
\maketitle

\noindent\emph{Source and licence.} \bf Equivariant critical point theory and bifurcation of $3d$ gravity-capillary Stokes waves. Version arXiv:2603.27847v2 (CC BY 4.0), distributed by arXiv under the Creative Commons Attribution 4.0 International licence, http://creativecommons.org/licenses/by/4.0/, as recorded in the arXiv licence metadata for this exact version and shown on the abstract page https://arxiv.org/abs/2603.27847v2. Full licence notice: \texttt{licenses/arxiv-2603.27847-CC-BY-4.0.txt}.\par\smallskip

\noindent\textit{Editorial scope.} Counted unit: the article's lemma cI (source0.tex lines 5803--5831). The article's complete original proof of it is delivered with it (source0.tex lines 5832--5981, one \texttt{proof} environment of 150 lines). No mathematics is written, repaired or re-derived: the statement and the proof are the article's own lines, in the article's own notation, and no retained line is substituted or re-flowed.  Retained uncounted as the source-local context the proof needs: lines 743--746; lines 844--847; lines 854--858; lines 2421--2426; lines 3787--3796; lines 3861--3866; lines 3880--3886; lines 3889--3891; lines 4372--4384; lines 4402--4427; lines 5661--5663; lines 5665--5666; lines 5680--5697.  Nothing was omitted from any retained span, so the package carries no omission record.  No retained line contains a citation, so the wrapper carries no bibliography and no bibliography entry.  The retained lines contain the article's own diagram code, which the wrapper typesets with the standard tikz packages; no image file is delivered and the package declares no asset dependency.  Editorial formatting only, in addition: an AMS \texttt{amsart} preamble that restates the article's own declarations and macros used by the retained lines, namely the environments \texttt{lemma}, \texttt{remark} on separate counters, together with the standard packages those lines need.  The retained environments print in this document as Lemma 1, Remark 1 in order of appearance, whereas the article prints the same items with the numbering of its own full document.  The complete native source of the article as distributed in the arXiv e-print for this exact version is bundled unchanged as \texttt{sources/197388/source0.tex}.
 \begin{equation}\label{Def cV}
        \cV:= \big\{
        j\in \Gamma'\setminus\{0\} \ | \ \omega(j)=c_*\cdot j \big\} 
    \end{equation}

\begin{equation}\label{defC}
    \cC := \Big\{ \sum_{j \in \cV } 
    \tau_j \, j  \ \big| \ \tau_j \leq 0 \Big\}
\end{equation}

\begin{equation}\label{eq:ilcono}
%\cC \subset \big\{ a \in \R^2 \ | \ a \cdot c_* \leq 0  \big\} \, ,
%\quad 
\cC \setminus \{0\} \subset \big\{ a \in \R^2 \ | \ a \cdot c_* < 0  \big\} \, . 
\end{equation} 

\begin{equation}
\label{eqrela}
(x,y,0)\sim(x,y',0) \, ,\
(x,y,1)\sim(x',y,1) \, ,\quad
\forall x,x'\in X \, ,\ y,y'\in Y \, .
\end{equation}

     \begin{lemma}\label{ortoges}
    {\bf (Characterization of $2d$-waves)}
    Any function 
            $u \in L^2  $ is a 
            $ 2 d$-wave along the direction $ [j] $, i.e.  $ u\in \cX_{[j]}^{\twod} $, if and only if 
    \begin{equation}\label{One dimensional iff simmetric}
         \tau_\theta u=u \, , \quad \forall \theta\in 
                {\mathtt T}_{[j]} \, , \quad \text{i.e.} \quad {\mathtt T}_{[j]}\subseteq (\T_\Gamma^2)_u\, .
            \end{equation}
    \end{lemma}

\begin{equation}\label{V2d}
{\bf V}^{\twod} := \bigcup_{\hat\jmath \in \cD} V_{\hat\jmath}
\, , \qquad 
V_{\hat\jmath} = \big\{ v \in V \ | \ 
\Pi_{\hat\jmath}^\bot v = 0  \big\} \, , 
\end{equation}

 \begin{remark}\label{rem:Vinv}
By Lemma \ref{ortoges}
an equivariant function $ f : V \to V $ 
maps each subspace $ V_{\hat\jmath} \to V_{\hat\jmath}  $,  for any 
$ \hat\jmath \in \cD  $,  and $ V \setminus {\bf V}^{\twod}
\to V \setminus {\bf V}^{\twod} $. 
\end{remark}

    \begin{equation}\label{eq:cI}
        \cI(v)=-\frac12\sum_{j\in \cV} \hat\jmath \|\pi_{j} v\|^2=-\frac12\sum_{\hat\jmath \in \cD} \hat\jmath \|\Pi_{\hat\jmath} v\|^2
    \end{equation}

    \begin{equation}\label{Def a}
        \cA (v) :=
        \begin{pmatrix}
            \cA_{11} (v)  & \cA_{12} (v)   \\
             \cA_{12} (v)  &  \cA_{22} (v) 
        \end{pmatrix}  = 
        \begin{pmatrix}
            \langle \nabla_v \cI_1(v), \nabla_v \cI_1(v)\rangle  &
            \langle \nabla_v \cI_2(v), \nabla_v \cI_1(v)\rangle \\
            \langle \nabla_v \cI_1(v), \nabla_v \cI_2(v)\rangle  &
            \langle \nabla_v \cI_2(v), \nabla_v \cI_2(v)\rangle 
        \end{pmatrix}  
    \end{equation} 

    \begin{lemma} \label{Properties of a}
   The matrix  $ \cA(v)$ in \eqref{Def a} is  equal to   
        \begin{equation} \label{NablaI(v)NablaI(v) in coordinates}
            \cA(v)= \sum_{j\in \cV}
            \|\pi_{j} v\|^2\, \hat\jmath\otimes \hat\jmath = \sum_{{\hat \jmath}\in \cD}
            \|\Pi_{\hat \jmath} v\|^2
            \, \hat\jmath\otimes \hat\jmath \, ,
            \quad  \forall 
            v \in V \, , 
        \end{equation}
    where $ \hat\jmath\otimes \hat\jmath $ is the symmetric, rank $ 1 $, positive semi-definite  matrix  $\hat \jmath \otimes \hat \jmath:=\hat \jmath \hat \jmath^\top$.
    For any $v\in V$
    \begin{equation}\label{detAespli}
     \det(\cA(v))=\frac12\sum_{\hat \jmath,\hat \jmath' \in \cD} \det([\hat \jmath| \hat \jmath'])^2\, \|\Pi_{\hat \jmath} v\|^2\|\Pi_{\hat \jmath'} v\|^2
    \end{equation}
    where $[\hat\jmath| \hat \jmath']$ is the matrix with columns  $\hat \jmath$ and $\hat \jmath'$.
      The matrix $ \cA(v) $ is singular    on the set 
        \begin{equation}\label{Def S*}
            \Big\{ v \in V \ |\ \det( \cA ( v))  =0 \Big\} = {\bf V}^{\twod}  = 
            \bigcup_{{\hat \jmath} \in \cD} V_{\hat \jmath} \, ,    
        \end{equation}
      % (cf. \eqref{V2d}), 
      more precisely 
    $ \det (\cA (v)) > 0 $  on  
$ d(v,{\bf V}^{\twod}) > 0 $. 
        \end{lemma}

\begin{equation}\label{eq:def V-0+}
    V = V_- \oplus V_0 \oplus V_+   
 \end{equation}   

\begin{equation}\label{eq:def cD-0+}
V_-:=\bigoplus_{\hat \jmath \in \cD_- }V_{\hat \jmath}\, , \quad  V_0:=\bigoplus_{\hat \jmath \in \cD_0} V_{\hat \jmath} \, , \quad  V_+:=\bigoplus_{\hat \jmath \in \cD_+ }V_{\hat \jmath}  \, , \end{equation}

\begin{remark}\label{lem:linese}
If $a\in \interior(\cC)$  is in the interior of
the  cone  defined   in \eqref{defC} 
  then  
  $\cD_\pm \not=\emptyset$ and so 
  $\dim(V_\pm ) \geq 2$.
 % Here  $\dim(\cdot)$ denotes the real dimension of a vector space. 
%Let $\cC$ be    
 %[$(\beth)$]  {\bf (Boundary case)}  
%    
If 
$a \in \partial \cC\setminus\{0\} $ 
 then 
  $\cD_0 \not =\emptyset $, so  
$\dim(V_0)\geq 2 $, and   
%since $\cC$ is  convex, 
either $\cD_- =\emptyset$ or $\cD_+ =\emptyset$.   
\end{remark}

        \begin{lemma}{\bf (Level sets of $\cI$)}\label{prop:levelsets of cI}
           If  $a\in \cC \setminus \{0\} $ then 
    \\[1mm]     
        $(\beth)$
               If $a\in \partial \cC\setminus\{0\}$ then  $\cI^{-1}(a) \subset V_0 $ is $\T^2_\Gamma$-equivariantly diffeomorphic to $S(V_0)$.
               \\[1mm]
  $(\daleth)$
                If $a\in \interior(\cC)\setminus\{0\}  $ then the following exclusive cases may occur:
                \begin{itemize}
                    \item[$(\daleth$1)] 
                     If $a\not \parallel \hat \jmath$ for any $\hat \jmath\in \cD$ then  $\cI^{-1}(a)$ is a closed submanifold of $V$  contained in $V\setminus 
                    {\bf V}^{\twod} $ which is 
                    $\T^2_\Gamma$-equivariantly diffeomorphic to $S(V_-)\times S(V_+) $.
                    
                    \item[$(\daleth$2)] 
                     If $a\parallel \hat \jmath_0$ for some resonant wave direction $\hat \jmath_0\in \cD$
                    then  
 $       \cI^{-1}(a)\setminus {\bf V}^{\twod}  = 
                     \cI^{-1}(a)\setminus  V_0  $
is a submanifold of $  V $, 
       $  \cI^{-1}(a)\cap {\bf V}^{\twod}  = \cI^{-1}(a)\cap V_0 $ is a submanifold of  $ V_0 = V_{\hat \jmath_0} $, diffeomorphic to 
       $ S(V_0) $ 
                      and there is a $\T^2_\Gamma$-equivariant homeomorphism 
                    \begin{equation}\label{gamma homeomorphism_2}
                        \gamma_1:\cI^{-1}(a)\toup^{\cong}[S(V_-)\times S(V_+)] \star S(V_0)\, ,
                        \quad \gamma_1(\cI^{-1}(a) {\cap {\bf V}^{2d}})=S(V_0) \, .
                    \end{equation}
                \end{itemize}
        \end{lemma}    

        \begin{proof} 
       
    Projecting  $ \cI (v) $ in \eqref{eq:cI}  
       along  $ a $ and 
       $ a^\bot $, 
    we write
    \begin{equation}
    \label{J-1a}
    \cI^{-1} (a) 
    = \Big\{ v\in V \ | \
    \sum_{\hat \jmath\in \cD}\frac{-\hat \jmath \cdot a}{2|a|^2}\|\Pi_{\hat \jmath}v\|^2=1 \, ,  \ \sum_{\hat \jmath\in \cD}\frac{-\hat \jmath \cdot a^\bot}{2|a|^2}\|\Pi_{\hat \jmath}v\|^2=0 
    \Big\} \, . 
\end{equation}
By \eqref{eq:ilcono}, 
we have $a\cdot c_*<0 $ and
 summing the second equation in \eqref{J-1a} multiplied by $\frac{ a^\bot\cdot c_*}{a\cdot c_*}$ to the first equation, 
 and decomposing  $  c_* = 
a (c_* \cdot a)|a|^{-2} + a^\bot 
(c_* \cdot a^\bot) |a|^{-2} $, we deduce that 
\begin{equation}\label{J-2a}
  \cI^{-1} (a)  = \Big\{ v\in V \ | \
    \cB(v)=1 \, ,  \ \cK(v)=0 
    \Big\}
    \end{equation}
    where
     \begin{equation}\label{cBcK}
            \cB(v):=\frac{1}{2|a\cdot c_*|}\sum_{\hat \jmath\in \cD} c_*\cdot \hat \jmath \, \|\Pi_{\hat \jmath} v\|^2 \, , \quad 
            \cK(v):=\sum_{\hat \jmath \in \cD} \hat \jmath \cdot a^\bot \|\Pi_{\hat \jmath} v\|^2\, .
    \end{equation}
    Note that $c_*\cdot \hat \jmath  >0$ for any $ \hat \jmath \in \cD$ (cf. \eqref{Def cV}) and hence $ \sqrt{\cB(v)}$ is a norm on $ V $. 

    We define the $\T^2_\Gamma$-equivariant diffeomorphism
\begin{equation}\label{diffeoxi}
\xi:V\setminus\{0\}\to V\setminus\{0\} \, , \quad
            \xi(v):= g(v) \psi(v)
        \, ,  \quad  
g(v) :=  \frac{\vvv v\vvv}{\sqrt{\cB( \psi(v))}}     \, , 
\end{equation}
where $ \psi:V\to V $
            is the linear isomorphism    
      \begin{equation}\label{defpsiso}
         \psi(v) := 
         \sum_{\hat \jmath\in \cD_-} |\hat \jmath \cdot a^\bot|^{-\frac12}\, \Pi_{\hat \jmath} v  +
        \sum_{\hat \jmath\in \cD_0} \Pi_{\hat \jmath} v
         + \sum_{\hat \jmath\in \cD_+} (\hat \jmath\cdot {a^\bot})^{-\frac12}\, \Pi_{\hat \jmath} v 
        \end{equation}
  (it is well defined in view of \eqref{eq:def cD-0+}) 
  and
\begin{equation}\label{def3bar}
\vvv v\vvv^2:=\tfrac 12\|\Pi_- v\|^2+\tfrac 12\|\Pi_+ v\|^2+\|\Pi_0 v\|^2 \, , 
\end{equation}
denoting 
$ \Pi_\pm :V\to V_\pm $ and $\Pi_0:V\to V_0$ 
 the orthogonal projectors associated to the decomposition \eqref{eq:def V-0+}.   
Note that $ g(v) $ in \eqref{diffeoxi} is a  
positive,  $ C^\infty (V \setminus \{0\}, \R) $
homogenous function  
of degree $ 0 $.
Under the diffeomorphism  $ \xi $ in \eqref{diffeoxi}
the functions $ \cB $ and $ \cK $ in \eqref{cBcK} transform into
\begin{equation}\label{futras}
            (\cB \circ \xi)(v) =\vvv v\vvv^2\, ,\quad 
            (\cK \circ \xi)(v) =g^2(v)\big(\|\Pi_+ v\|^2 -\|\Pi_- v\|^2\big)\, ,
    \quad \forall 
    v\in V\setminus\{0\} \,            .
\end{equation}
As a consequence of \eqref{futras} and \eqref{def3bar}, the map   
$\xi$ defines a 
        $ \T^2_\Gamma $-equivariant diffeomorphism which maps the set $\cI^{-1}(a)$ 
        in \eqref{J-2a}  onto 
\begin{equation}\label{eq:Sigma}
        \Sigma := \Big\{ v \in V 
        \ | \  
            \tfrac12\|\Pi_- v\|^2+\tfrac 12\|\Pi_+ v\|^2+\|\Pi_0 v\|^2=1 \, , \|\Pi_+ v\|^2 =\|\Pi_- v\|^2 \Big\}  \, . 
        \end{equation} 
       $(\beth)$ {\sc Case $a\in \partial \cC\setminus\{0\}$}.
        In view of  Remark \ref{lem:linese},
        either $ \cD_- $ or $ \cD_+ $ is empty, say $ \cD_- = \emptyset $. Then   
        the set $ \Sigma $ in \eqref{eq:Sigma} 
        reduces to
        $\Sigma= \{ v \in V \ | \  \Pi_+ v = 0 \, , \ 
        \| \Pi_0 v \|^2 = 1 \} \equiv S(V_0) $ and  
        $\xi:
           S(V_0) \to \cI^{-1}(a) 
           = \xi  (S(V_0)) \subset V_0 $ 
           (the map $ \psi$ in \eqref{defpsiso}, thus $ \xi $,  maps $  V_0\setminus\{0\} \to V_0\setminus\{0\} $) is a $\T^2_\Gamma$-equivariant diffeomorphism.
           \\[1mm]
     $(\daleth)$ {\sc Case  $a\in \interior(\cC)$}.
        Remark \ref{lem:linese} shows that 
        $\dim(V_\pm )  \geq 2$. 
        We distinguish two  sub-cases:
        \begin{itemize}
            \item[$(\daleth$1)] Assume $ a \not \parallel \hat \jmath $ for any $\hat \jmath\in \cD$. Then $\cD_0=\emptyset$ and 
           \eqref{eq:Sigma} reduces to
        $\Sigma= \{ v \in V \ | \  
        \| \Pi_- v \|^2 = 
        \| \Pi_+ v \|^2 = 1 \} 
           $
           $ \equiv S(V_-)\times S(V_+)\subset V\setminus{\bf V}^{\twod}$. Hence $\xi: S(V_-) \times S(V_+) \to \cI^{-1}(a) $ 
             is a $\T^2_\Gamma$-equivariant diffeomorphism and Remark \ref{rem:Vinv}  implies  $ \cI^{-1}(a)\subset V\setminus{\bf V}^{\twod}$.
            \item[$(\daleth$2)] Assume $a  \parallel \hat \jmath_0 $ for some $\hat \jmath_0 \in \cD $.
Since the diffeomorphism 
$ \xi $ in \eqref{diffeoxi} maps each 
$ V_{\hat \jmath} \setminus \{0  \} $ in itself and  
$ V \setminus  {\bf V}^{\twod} $ in itself, we have
$$
\xi : 
\underbrace{\Sigma \cap 
{\bf V}^{\twod}}_{= S(V_0) \, \text{by} \,  \eqref{eq:Sigma}} \to \underbrace{\cI^{-1}(a)\cap {\bf V}^{\twod}}_{= \cI^{-1}(a)\cap V_0} 
\, , \quad \xi : \Sigma \setminus  
{\bf V}^{\twod} \to \cI^{-1}(a)\setminus  {\bf V}^{\twod} =
\cI^{-1}(a)\setminus V_0\, , 
% = \xi (\Sigma \cap {\bf V}^{\twod}) 
$$         
and  $ \cI^{-1}(a)\setminus {\bf V}^{\twod}$  is a sub-manifold of $V$ of codimension two because, 
for any $v\in V\setminus {\bf V}^{\twod}$ the vectors  $\{\nabla_v \cI_1(v), \nabla_v \cI_2 (v)\}$ are linearly independent 
by Lemma \ref{Properties of a}.        
         We now prove that 
         $ \Sigma $
         is homeomorphic to the join 
         topological space $ M \star Y $
         where $ M := S(V_-)\times S(V_+) $ and 
         $ Y := S(V_0) $. 
        The map
        $$
        f : M \times   Y \times [0,1] \to \Sigma \, , \quad  (v_1, v_0,t) \mapsto 
        f(v_1, v_0, t) := \sqrt{1-t^2}\, v_1 + {t}\, v_0 \, ,
        $$
        is 
        continuous, surjective, $ \T^2_\Gamma $-equivariant,  
and defines a map $ \tilde f $ on the quotient space
\[
\begin{tikzcd}[row sep=large]
 M \times Y \times [0,1] 
 \arrow[d,"\pi"']  \arrow[r,"f"]
 & 
\Sigma 
 \\
M \star Y = \frac{M \times Y\times[0,1]}\sim 
 \arrow[ur,"\tilde f "',"\cong"]
 & 
\end{tikzcd}
\]
where $ \sim $ is the equivalence relation defined in \eqref{eqrela}.
The map $ \tilde f $ is continuous, bijective and since $ M, Y $ are compact it is a $ \T^2_\Gamma $-equivariant homeomorphism between 
$ M \star Y $ and $ \Sigma $. The homeomorphism \eqref{gamma homeomorphism_2} is 
$ \gamma_1 := \tilde f^{-1} \circ \xi^{-1} $ (note that $ \tilde f (Y) = S(V_0) $).        
        \qedhere
        \end{itemize}
        \end{proof}

\end{document}
